% !PRINTED 503
\weberpage{503}
%<<CONT
schwindet, sondern beliebig viele, da man jedes $\Phi(t)$ mit einer
beliebigen Function der gleichen Form multipliciren kann. Auch
kann es vorkommen, dass schon ein Product von weniger als
$n$ Factoren von $N(t-\omega)$ rationale Co\"efficienten erhält, woraus
Functionen $\Phi(t)$ von niedrigerem als dem $n^{\text{ten}}$ Grade entspringen,
die für $t = \omega$ verschwinden. Daraus folgt:

7. \gesperrt{Es giebt für jedes Functional $\omega$ des Körpers $\Omega$
unendlich viele Functionen $\Phi(t)$, die für $t = \omega$
verschwinden, in denen die höchste Potenz von $t$
den Co\"efficienten 1 hat, und deren übrige Co\"effi\-
cienten in $\overline{R}$ enthalten sind.}

Wir sagen dann, $\omega$ ist eine Wurzel der Gleichung
\[
\Phi(t) = 0.
\]

Unter den verschiedenen Functionen $\Phi(t)$, deren Existenz
im Satze 7. ausgesprochen ist, giebt es eine und nur eine $F(t)$
von möglichst niedrigem Grade. Denn existiren zwei solche
Functionen $F(t)$ und $F_1(t)$ von gleichem Grade $m$, so ist die
Differenz $F(t) - F_1(t)$ in Bezug auf $t$ höchstens vom Grade
$m - 1$ und verschwindet für $t = \omega$. Dividirt man durch den
Co\"efficienten der höchsten Potenz von $t$, so erhält man eine
Function $\Phi(t)$ von niedrigerem Grade als $F(t)$, die nach der
Voraussetzung über $F(t)$ nicht existirt.

Die Function $F(t)$ kann im Körper $\overline{R}$ nicht in Factoren
zerlegt werden, die ganze Functionen von $t$ sind. Denn zerfiele
sie in mehrere Factoren derselben Form $F_1(t)$, $F_2(t)$, so müsste
einer dieser Factoren, die doch alle von niedrigerem Grade als
$F(t)$ selbst sind, für $t = \omega$ verschwinden, entgegen unserer
Voraussetzung.

Jede Function $\Phi(t)$ des Satzes 7. ist dann durch diese
irreducible Function $F(t)$ theilbar, so dass der Quotient
$\Phi(t):F(t)$ eine ganze Function von $t$ in $\overline{R}$ ist, in der die
höchste Potenz von $t$ den Co\"efficienten 1 hat (vgl. Bd. I,
\S.\,141, I.).

Unter den Functionen $\Phi(t)$ findet sich, wie schon bemerkt,
auch die Norm $N(t-\omega)$, und folglich ist $N(t-\omega)$ durch $F(t)$
theilbar. Sind beide Functionen von gleichem Grade, so ist
$N(t-\omega) = F(t)$. Ist aber der Grad von $F(t)$ niedriger als $n$,
so verschwindet der Quotient $N(t-\omega):F(t)$ wenigstens noch
für eines der mit $\omega$ conjugirten Functionale, und folglich für
alle; folglich auch für $t = \omega$, und daher ist dieser Quotient
%CONT>>
